\textbf{Phần 1. Trắc nghiệm lựa chọn}
\setcounterpageref{ex}{0}

\begin{ex} %[5X12V74]
Một trường tiểu học có $600$ học sinh. Biểu đồ quạt tròn cho biết cơ cấu học sinh tham gia câu lạc bộ thể thao gồm: Bóng đá chiếm $45\%$, Cầu lông chiếm $30\%$, Bơi lội chiếm $25\%$. Số học sinh tham gia câu lạc bộ bóng đá nhiều hơn câu lạc bộ bơi lội là
\choice
{$90$ em}
{\True $120$ em}
{$150$ em}
{$60$ em}
\loigiai{
Tỉ số phần trăm học sinh thích bóng đá nhiều hơn bơi lội là:
$$45\% - 25\% = 20\%$$
Số học sinh tham gia câu lạc bộ bóng đá nhiều hơn bơi lội là:
$$600 \times 20 : 100 = 120\text{ (em)}$$
\textbf{Đáp án của câu hỏi là B}
}
\end{ex}

\begin{ex} %[5X12V74]
Biểu đồ quạt tròn biểu diễn thể loại sách yêu thích trong một thư viện gồm: Truyện tranh $40\%$, Cổ tích $35\%$, Khoa học $25\%$. Biết số bạn thích truyện tranh nhiều hơn số bạn thích truyện khoa học là $45$ bạn. Tổng số bạn tham gia cuộc khảo sát là
\choice
{$250$ bạn}
{\True $300$ bạn}
{$350$ bạn}
{$400$ bạn}
\loigiai{
Hiệu tỉ số phần trăm giữa truyện tranh và khoa học là:
$$40\% - 25\% = 15\%$$
Tổng số bạn tham gia khảo sát là:
$$45 : 15 \times 100 = 300\text{ (bạn)}$$
\textbf{Đáp án của câu hỏi là B}
}
\end{ex}

\begin{ex} %[5X12V74]
Trong hộp có các viên bi màu đỏ và màu xanh. Một học sinh thực hiện lấy ngẫu nhiên một viên bi, ghi lại màu rồi trả lại vào hộp. Sau $80$ lần lấy bi liên tiếp, có $48$ lần lấy được viên bi màu đỏ. Tỉ số phần trăm số lần lấy được viên bi màu xanh là
\choice
{$35\%$}
{\True $40\%$}
{$45\%$}
{$60\%$}
\loigiai{
Số lần lấy được viên bi màu xanh là:
$$80 - 48 = 32\text{ (lần)}$$
Tỉ số phần trăm số lần lấy được bi màu xanh là:
$$32 : 80 = 0{,}4 = 40\%$$
\textbf{Đáp án của câu hỏi là B}
}
\end{ex}

\begin{ex} %[5X12V74]
Điểm kiểm tra môn Toán của $4$ bạn học sinh trong một tổ là $7, 8, 8, 9$. Sau khi bổ sung thêm điểm kiểm tra của bạn học sinh thứ năm thì điểm trung bình của cả $5$ bạn là $8{,}2$ điểm. Điểm kiểm tra của bạn thứ năm là
\choice
{$8$ điểm}
{\True $9$ điểm}
{$10$ điểm}
{$8{,}5$ điểm}
\loigiai{
Tổng điểm của 4 bạn ban đầu là:
$$7 + 8 + 8 + 9 = 32\text{ (điểm)}$$
Tổng điểm của cả 5 bạn là:
$$8{,}2 \times 5 = 41\text{ (điểm)}$$
Điểm kiểm tra của bạn học sinh thứ năm là:
$$41 - 32 = 9\text{ (điểm)}$$
\textbf{Đáp án của câu hỏi là B}
}
\end{ex}

\begin{ex} %[5X12V74]
Một trang trại nuôi gia cầm gồm: Gà ($50\%$), Vịt ($30\%$), Ngan ($12\%$) và Ngỗng. Biết trang trại có tất cả $2\,500$ con gia cầm. Số con ngỗng trang trại đang nuôi là
\choice
{$150$ con}
{\True $200$ con}
{$250$ con}
{$300$ con}
\loigiai{
Tỉ số phần trăm số ngỗng trong trang trại là:
$$100\% - (50\% + 30\% + 12\%) = 8\%$$
Số con ngỗng trang trại nuôi là:
$$2\,500 \times 8 : 100 = 200\text{ (con)}$$
\textbf{Đáp án của câu hỏi là B}
}
\end{ex}

\begin{ex} %[5X12V74]
Gieo một con xúc xắc $120$ lần, quan sát thấy số lần xuất hiện mặt có số chấm chia hết cho 2 (mặt chẵn: 2, 4, 6) là $66$ lần. Tỉ số của số lần xuất hiện mặt có số chấm lẻ so với tổng số lần gieo là
\choice
{$\dfrac{11}{20}$}
{\True $\dfrac{9}{20}$}
{$\dfrac{9}{11}$}
{$\dfrac{1}{2}$}
\loigiai{
Số lần xuất hiện mặt có số chấm lẻ là:
$$120 - 66 = 54\text{ (lần)}$$
Tỉ số giữa số lần xuất hiện mặt lẻ và tổng số lần gieo là:
$$\dfrac{54}{120} = \dfrac{9}{20}$$
\textbf{Đáp án của câu hỏi là B}
}
\end{ex}

\begin{ex} %[5X12V74]
Một cửa hàng bán được $4$ loại bánh trong ngày: Bánh dâu ($35$ chiếc), Bánh cam ($45$ chiếc), Bánh sô-cô-la ($80$ chiếc), Bánh phô mai ($40$ chiếc). Tỉ số phần trăm số bánh sô-cô-la so với tổng số bánh đã bán là
\choice
{$35\%$}
{\True $40\%$}
{$45\%$}
{$50\%$}
\loigiai{
Tổng số bánh cửa hàng đã bán là:
$$35 + 45 + 80 + 40 = 200\text{ (chiếc)}$$
Tỉ số phần trăm số bánh sô-cô-la là:
$$80 : 200 = 0{,}4 = 40\%$$
\textbf{Đáp án của câu hỏi là B}
}
\end{ex}

\begin{ex} %[5X12V74]
Một nhà máy kiểm tra chất lượng $500$ linh kiện điện tử thì phát hiện có $15$ linh kiện bị lỗi. Tỉ lệ linh kiện đạt chuẩn chất lượng của lô hàng đó là
\choice
{$96\%$}
{\True $97\%$}
{$98\%$}
{$95\%$}
\loigiai{
Số linh kiện đạt chuẩn chất lượng là:
$$500 - 15 = 485\text{ (linh kiện)}$$
Tỉ số phần trăm linh kiện đạt chuẩn là:
$$485 : 500 = 0{,}97 = 97\%$$
\textbf{Đáp án của câu hỏi là B}
}
\end{ex}

\begin{ex} %[5X12V74]
Một gia đình có tổng thu nhập mỗi tháng là $24$ triệu đồng. Biểu đồ quạt tròn cho biết cơ cấu chi tiêu: Tiền ăn uống chiếm $45\%$, Tiền học cho các con chiếm $25\%$, Chi phí sinh hoạt khác chiếm $20\%$, còn lại là Tiền tiết kiệm. Số tiền gia đình đó tiết kiệm được mỗi tháng là
\choice
{$1{,}8\text{ triệu đồng}$}
{\True $2{,}4\text{ triệu đồng}$}
{$3\text{ triệu đồng}$}
{$3{,}6\text{ triệu đồng}$}
\loigiai{
Tỉ số phần trăm số tiền tiết kiệm mỗi tháng là:
$$100\% - (45\% + 25\% + 20\%) = 10\%$$
Số tiền tiết kiệm được mỗi tháng là:
$$24 \times 10 : 100 = 2{,}4\text{ (triệu đồng)}$$
\textbf{Đáp án của câu hỏi là B}
}
\end{ex}

\begin{ex} %[5X12V74]
Trong một trò chơi bốc thăm trúng thưởng có $100$ lượt bốc thăm, ghi nhận có $28$ lượt trúng giải thưởng. Tỉ số phần trăm số lượt không trúng giải thưởng là
\choice
{$68\%$}
{\True $72\%$}
{$75\%$}
{$82\%$}
\loigiai{
Số lượt không trúng giải thưởng là:
$$100 - 28 = 72\text{ (lượt)}$$
Tỉ số phần trăm số lượt không trúng giải là:
$$72 : 100 = 72\%$$
\textbf{Đáp án của câu hỏi là B}
}
\end{ex}

\begin{ex} %[5X12V74]
Lượng mưa đo được trong $4$ tháng mùa hè tại một địa phương lần lượt là: Tháng 5 ($180\text{ mm}$), Tháng 6 ($240\text{ mm}$), Tháng 7 ($310\text{ mm}$), Tháng 8 ($270\text{ mm}$). Lượng mưa trung bình mỗi tháng trong 4 tháng mùa hè đó là
\choice
{$240\text{ mm}$}
{\True $250\text{ mm}$}
{$260\text{ mm}$}
{$275\text{ mm}$}
\loigiai{
Tổng lượng mưa của 4 tháng là:
$$180 + 240 + 310 + 270 = 1\,000\text{ (mm)}$$
Lượng mưa trung bình mỗi tháng là:
$$1\,000 : 4 = 250\text{ (mm)}$$
\textbf{Đáp án của câu hỏi là B}
}
\end{ex}

\begin{ex} %[5X12V74]
Trong biểu đồ hình quạt tròn, phần quạt biểu diễn số học sinh đạt giải quốc tế chiếm $30\%$ diện tích toàn bộ hình tròn. Góc ở tâm tương ứng với phần hình quạt đó có số đo là
\choice
{$90^\circ$}
{\True $108^\circ$}
{$120^\circ$}
{$72^\circ$}
\loigiai{
Góc ở tâm của cả hình tròn là $360^\circ$.\\
Góc ở tâm của phần quạt chiếm $30\%$ là:
$$360^\circ \times 30 : 100 = 108^\circ$$
\textbf{Đáp án của câu hỏi là B}
}
\end{ex}

\textbf{Phần 2. Trắc nghiệm đúng sai}
\setcounterpageref{ex}{0}

\begin{ex} %[5X12V74]
Khảo sát phương tiện đi lại đến trường của $200$ học sinh khối 5: Đi xe đạp ($90$ em), Đi bộ ($50$ em), Đi xe buýt ($40$ em), Được bố mẹ chở bằng xe máy ($20$ em). Xét tính đúng sai của các khẳng định sau:
\choiceTF
{\True Số học sinh đi xe đạp chiếm $45\%$ tổng số học sinh.}
{\True Số học sinh đi bộ chiếm $25\%$ tổng số học sinh.}
{Số học sinh đi xe buýt chiếm $15\%$ tổng số học sinh.}
{\True Số học sinh được bố mẹ chở chiếm $10\%$ tổng số học sinh.}
\loigiai{
\begin{itemchoice}
    \itemch Đúng. Xe đạp: $90 : 200 = 45\%$.
    \itemch Đúng. Đi bộ: $50 : 200 = 25\%$.
    \itemch Sai. Xe buýt: $40 : 200 = 20\% \neq 15\%$.
    \itemch Đúng. Bố mẹ chở: $20 : 200 = 10\%$.
\end{itemchoice}
}
\end{ex}

\begin{ex} %[5X12V74]
Bạn Linh thực hiện tung một đồng xu liên tiếp $150$ lần và ghi lại kết quả: Mặt Sấp xuất hiện $84$ lần, mặt Ngửa xuất hiện $66$ lần. Xét tính đúng sai của các khẳng định sau:
\choiceTF
{\True Tỉ số của số lần xuất hiện mặt Sấp so với tổng số lần tung là $\dfrac{14}{25}$.}
{\True Tỉ số phần trăm số lần xuất hiện mặt Ngửa là $44\%$.}
{\True Mặt Sấp xuất hiện nhiều hơn mặt Ngửa là $18$ lần.}
{Tỉ số phần trăm số lần xuất hiện mặt Sấp là $55\%$.}
\loigiai{
\begin{itemchoice}
    \itemch Đúng. Tỉ số mặt Sấp: $\dfrac{84}{150} = \dfrac{14}{25}$.
    \itemch Đúng. Phần trăm mặt Ngửa: $66 : 150 = 0{,}44 = 44\%$.
    \itemch Đúng. Chênh lệch: $84 - 66 = 18$ lần.
    \itemch Sai. Phần trăm mặt Sấp: $84 : 150 = 0{,}56 = 56\% \neq 55\%$.
\end{itemchoice}
}
\end{ex}

\textbf{Phần 3. Trả lời ngắn}
\setcounterpageref{ex}{0}

\begin{ex} %[5X12V74]
Gieo một con xúc xắc $80$ lần, thấy mặt $6$ chấm xuất hiện $16$ lần. Tính tỉ số phần trăm số lần xuất hiện mặt $6$ chấm. (Chỉ ghi số).
\shortans{20}
\loigiai{
Tỉ số phần trăm số lần xuất hiện mặt $6$ chấm là:
$$16 : 80 = 0{,}2 = 20\%$$
\textbf{Đáp số:} $20$
}
\end{ex}

\begin{ex} %[5X12V74]
Biểu đồ quạt tròn cho biết học sinh đạt loại Giỏi chiếm $40\%$, loại Khá chiếm $45\%$, còn lại là loại Trung bình. Biết khối 5 có tất cả $300$ học sinh, tính số học sinh đạt loại Trung bình.
\shortans{45}
\loigiai{
Tỉ số phần trăm học sinh đạt loại Trung bình là:
$$100\% - (40\% + 45\%) = 15\%$$
Số học sinh đạt loại Trung bình là:
$$300 \times 15 : 100 = 45\text{ (học sinh)}$$
\textbf{Đáp số:} $45$
}
\end{ex}

\begin{ex} %[5X12V74]
Điểm kiểm tra môn Tiếng Anh của $5$ bạn trong một nhóm lần lượt là: $8, 8, 9, 9, 10$. Khi có thêm bạn thứ sáu tham gia nhóm thì điểm trung bình của cả $6$ bạn là $9$ điểm. Tính điểm kiểm tra của bạn thứ sáu.
\shortans{10}
\loigiai{
Tổng điểm của 5 bạn ban đầu là:
$$8 + 8 + 9 + 9 + 10 = 44\text{ (điểm)}$$
Tổng điểm của cả 6 bạn là:
$$9 \times 6 = 54\text{ (điểm)}$$
Điểm của bạn thứ sáu là:
$$54 - 44 = 10\text{ (điểm)}$$
\textbf{Đáp số:} $10$
}
\end{ex}

\begin{ex} %[5X12V74]
Một cửa hàng bán được $500$ gói trà gồm $3$ loại: Trà xanh chiếm $50\%$, Trà ô long chiếm $32\%$, còn lại là Trà hoa cúc. Tính số gói Trà hoa cúc đã bán.
\shortans{90}
\loigiai{
Tỉ số phần trăm số gói Trà hoa cúc là:
$$100\% - (50\% + 32\%) = 18\%$$
Số gói Trà hoa cúc đã bán là:
$$500 \times 18 : 100 = 90\text{ (gói)}$$
\textbf{Đáp số:} $90$
}
\end{ex}

\begin{ex} %[5X12V74]
Quay một vòng quay may mắn $60$ lần, kim dừng lại ở ô trúng thưởng $15$ lần. Tính tỉ số phần trăm số lần không trúng thưởng. (Chỉ ghi số).
\shortans{75}
\loigiai{
Số lần quay không trúng thưởng là:
$$60 - 15 = 45\text{ (lần)}$$
Tỉ số phần trăm số lần không trúng thưởng là:
$$45 : 60 = 0{,}75 = 75\%$$
\textbf{Đáp số:} $75$
}
\end{ex}

\begin{ex} %[5X12V74]
Biểu đồ hình quạt tròn biểu diễn kết quả bầu cử lớp trưởng cho ba bạn: Bạn An ($45\%$), bạn Bình ($35\%$), bạn Cường ($20\%$). Biết bạn An nhận được nhiều hơn bạn Bình là $12$ phiếu bầu. Tính tổng số phiếu bầu của cả lớp.
\shortans{120}
\loigiai{
Hiệu tỉ số phần trăm giữa phiếu của An và Bình là:
$$45\% - 35\% = 10\%$$
Tổng số phiếu bầu của cả lớp là:
$$12 : 10 \times 100 = 120\text{ (phiếu)}$$
\textbf{Đáp số:} $120$
}
\end{ex}
